\uuid{tFK9}
\exo7id{2931}
\titre{exo7 2931}
\auteur{quercia}
\organisation{exo7}
\datecreate{2010-03-08}
\isIndication{false}
\isCorrection{true}
\chapitre{Nombres complexes}
\sousChapitre{Géométrie}
\cpge{mpsi-mp2i.nombres-complexes.c, pcsi.nombres-complexes.c, ptsi.nombres-complexes.c, mpsi-mp2i.nombres-complexes.e, pcsi.nombres-complexes.e, ptsi.nombres-complexes.e, mpsi-mp2i.polynomes-fractions-rationnelles.c, pcsi.polynomes.c, ptsi.polynomes.c, bcpst1.nombres-complexes.b}
\module{Algèbre}
\niveau{L1}
\difficulte{}

\contenu{
\question{
Trouver $a,b,c \in \mathbb{U}$ tels que $\begin{cases} a+b+c = 1\cr abc = 1.\end{cases}$
}
\reponse{
$(0,a,a+b,a+b+c = 1)$ forme un losange donc l'un des nombres vaut $1$
         et les deux autres sont oppos{\'e}s $ \Rightarrow  \{a,b,c\} = \{1,i,-i\}$.
}
}
